\chapter{La conduite du projet}
\label{ch:conduite}

\questionchapitre{Ces développements n'ont été ni commandés ni hérités d'un dispositif
existant. Comment sont-ils devenus le travail du programme ?}

\section{Les interlocuteurs}

Les développements présentés dans ce rapport ont été conçus et réalisés dans leur intégralité au
cours du postdoctorat, dans un dialogue continu avec trois interlocuteurs aux apports distincts.

\paragraph{Supervision scientifique : Nicolas Labroche (LIFAT).} Il a encadré l'orientation
scientifique des travaux, en particulier l'arbitrage
initial : accepter que le sujet d'explicabilité impose d'abord la construction de son terrain
expérimental, plutôt que de le traiter sur un jeu de données de circonstance. Cet arbitrage
détermine tout ce qui précède.

\paragraph{Collaboration technique : Hugo Breuillard (BRGM).} Data scientist en charge du projet
côté BRGM, il constitue le point de contact sur les questions de chaîne de données et de
stratégie de prévision. Les échanges ont porté sur l'articulation entre l'entrepôt développé ici
et les besoins de modélisation du programme, et sur la cohérence des indicateurs produits avec
ceux employés par ailleurs.

\paragraph{Expertise métier : Augustin Thomas (BRGM).} Hydrogéologue, il a été l'interlocuteur
sur l'ensemble des questions de domaine. C'est de ces échanges que procèdent les décisions qui
ancrent la plateforme dans la pratique hydrogéologique plutôt que dans une transposition
générique de méthodes d'apprentissage.

\section{Porter le besoin jusqu'à son adoption}

Une part du travail n'apparaît ni dans les dépôts ni dans les mesures. Ces développements
n'ont pas été commandés : je les ai proposés, puis portés jusqu'à leur adoption.

\paragraph{Identifier le besoin et le faire reconnaître.} J'aurais pu contourner l'absence de jeu
de données par une extraction de circonstance, suffisante pour produire un résultat. J'ai pris le
parti inverse : formuler le manque, le documenter, puis le soumettre aux partenaires pour savoir
s'ils le partageaient. Le recadrage mené avec le BRGM au début de l'année 2026, avec les
interlocuteurs techniques comme avec les experts métier, a tranché. Le manque n'était pas propre
à ce postdoctorat : le besoin d'un entrepôt consolidé, d'outils d'analyse et d'un environnement
d'entraînement de modèles était partagé par les équipes du BRGM comme par d'autres projets du
programme, chacun butant sur le même obstacle de son côté. C'est cette vérification, et elle
seule, qui a rendu légitime l'investissement de plusieurs mois que représente un entrepôt de
données.

\paragraph{Aller chercher les interlocuteurs.} Aucun de ces contacts n'était hérité d'un
dispositif existant. Je les ai pris directement, auprès des interlocuteurs techniques comme des
experts métier, en préparant chaque échange autour des questions dont la réponse conditionnait un
choix de conception. La section~\ref{sec:apport-metier} en montre le résultat dans le code.

\paragraph{Tenir le suivi dans la durée.} Un échange ponctuel ne produit pas d'effet. J'ai
maintenu le suivi avec le BRGM tout au long de la période, et c'est ce qui explique que les
décisions métier ne soient pas restées au stade de la discussion.

\paragraph{Ouvrir le socle au-delà de son périmètre.} J'ai enfin présenté ces travaux aux équipes
du projet ANR \textsc{Aida}, dont le domaine d'expérimentation recouvre en partie celui traité
ici. L'intention n'était pas de revendiquer un périmètre, mais l'inverse : signaler l'existence
d'une infrastructure disponible, pour que ceux qui en auraient l'usage n'aient pas à la
reconstruire. Une infrastructure dont personne ne connaît l'existence ne vaut pas mieux qu'une
infrastructure inexistante.

\begin{center}
\begin{tabularx}{\textwidth}{@{}lX@{}}
\toprule
\textbf{Nature du travail} & \textbf{Ce qu'il a produit} \\
\midrule
Identification et formulation du besoin & Un constat vérifié auprès des partenaires, et non une
hypothèse personnelle \\
Prise de contact et préparation des échanges & Des décisions de conception fondées sur
l'expertise du domaine plutôt que sur des conventions d'informaticien \\
Suivi dans la durée & Des orientations métier effectivement traduites dans le code \\
Ouverture à d'autres équipes & Un socle connu au-delà de son projet d'origine \\
\bottomrule
\end{tabularx}
\end{center}

Un entrepôt de données mutualisé n'est pas un objet technique : c'est un objet d'organisation.
Sa valeur tient à ce que plusieurs équipes acceptent de s'y référer, ce qui relève de la
conviction plus que de la programmation. Le chapitre~\ref{ch:socle} montre que les obstacles
restants à son élargissement sont eux aussi de cette nature.
